\SourceLeaf{42}

\begin{center}
\textbf{POSTHUMOUS PAPERS.}
\end{center}

\begin{center}
\textbf{ON THE INTEGRATION OF THE DIFFERENTIAL FORMULA}\\
$(1+n\cos\varphi)^{\nu}. d\varphi.$
\end{center}

\noindent\hspace*{\parindent}\textls{Preface.} The greater part of the calculus known as integral calculus is carried out by expanding and working with series. The only method you will find in it apart from trial and recalling known differentiations is to expand a given differential quantity into a series whose terms can each be integrated with ease. With a finite series one thereby arrives at an expression complete in every respect. An infinite series represents the sought relation between the unknowns less perfectly. Provided it proceeds in such a way as to permit approximation, which suitable treatment will achieve in most cases, it is rightly considered sufficient for ordinary purposes and equivalent to an equation, if indeed there is one between the given variables. Yet because the terms entering the series are usually lengthy and apparently irregular, it is certainly an arduous task to specify their law by a general expression and to establish that expression by a proof. Both are demanded by the imminent danger of error creeping into a calculation based on mere conjecture made to appear plausible, and by the dignity of science. The present short treatise will furnish an example of such an investigation. It deals with carrying out the integration of the formula $d\varphi(1+n\cos\varphi)^{\nu}$ by series, where $n$ and $\nu$ are specified numbers taken at will. This formula is of outstanding use in physical astronomy and seemed so important to the distinguished \textsc{Euler} that an entire chapter of his Institutionum Calculi Integralis\HistoricalNote{*)}{*)}{Ch. VI, Vol. I.} was unhesitatingly%
